\documentclass[10pt,a4paper]{amsart}
\usepackage[margin=19mm]{geometry}
\usepackage{amsmath,amssymb,mathtools}
\usepackage[scr=rsfs,cal=euler]{mathalfa}
\usepackage{xfrac}
\usepackage[unicode,hidelinks,bookmarks=false]{hyperref}
\newcommand{\ie}{i.e.\ }
\newcommand{\cf}{cf.\ }
\newcommand{\resp}{respectively\ }
\newcommand{\Subsection}{\subsection}
\providecommand{\nobreakdash}{\nobreak\hbox{-}}
\setlength{\emergencystretch}{2em}
\multlinegap0pt
\usepackage{amssymb}
\newtheorem{thm}{Theorem}
\newtheorem{lem}{Lemma}
\title{Hausdorff dimensions and quasisymmetric minimalities of some homogeneous Moran sets}
\author{Jun Li, Yanzhe Li, Pingping Liu}
\date{Source: arXiv:2510.00540v3 (CC BY 4.0)}
\begin{document}
\maketitle

\noindent\emph{Source and licence.} Hausdorff dimensions and quasisymmetric minimalities of some homogeneous Moran sets. Version arXiv:2510.00540v3 (CC BY 4.0), distributed by arXiv under the Creative Commons Attribution 4.0 International licence, http://creativecommons.org/licenses/by/4.0/, as recorded in the arXiv licence metadata for this exact version and shown on the abstract page https://arxiv.org/abs/2510.00540v3. Full licence notice: \texttt{licenses/arxiv-2510.00540-CC-BY-4.0.txt}.\par\smallskip

\noindent\textit{Editorial scope.} Counted unit: the article's lem l10 (source0.tex lines 882--895). The article's complete original proof of it is delivered with it (source0.tex lines 898--1137, one \texttt{proof} environment of 240 lines). No mathematics is written, repaired or re-derived: the statement and the proof are the article's own lines, in the article's own notation, and no retained line is substituted or re-flowed.  Retained uncounted as the source-local context the proof needs: lines 215--238; lines 250--270; lines 551--561; lines 723--739; lines 832--834; lines 855--865.  Nothing was omitted from any retained span, so the package carries no omission record.  No retained line contains a citation, so the wrapper carries no bibliography and no bibliography entry.  No retained line contains a figure environment, an image-inclusion command or drawing code, so no image is delivered and the package declares no asset dependency.  Editorial formatting only, in addition: an AMS \texttt{amsart} preamble that restates the article's own declarations and macros used by the retained lines, namely the environments \texttt{thm}, \texttt{lem} on separate counters, together with the standard packages those lines need.  The retained environments print in this document as Theorem 1, Lemma 1 in order of appearance, whereas the article prints the same items with the numbering of its own full document.  The complete native source of the article as distributed in the arXiv e-print for this exact version is bundled unchanged as \texttt{sources/186145/source0.tex}.
\begin{thm}\label{thm}
    Let
	$E\in \mathcal{M}(I_{0},\left\{n_{k}\right\},\left\{c_{k}\right\})$ be a homogeneous Moran set which satisfies the following condition:
there exist two sequences of nonnegative real numbers $\{L_k\}_{k\ge1}$ and $\{R_k\}_{k\ge 1}$, such that
\begin{equation}
 \begin{aligned}
    \eta_{\sigma,0}=L_{k+1},\quad
    \eta_{\sigma,n_{k+1}}=R_{k+1} \nonumber
 \end{aligned}
 \end{equation}
 for any $k\ge 0$, $\sigma\in D_{k}$.

And if for any $k\ge 1$, at least one of the following three conditions is satisfied:
\begin{itemize}
  \item[(A)]there exists $\omega_1>0$, such that $\bar\alpha_k\le \omega_1 \underline\alpha_k$;
  \item[(B)]there exists $\omega_2>0$, such that $\bar\alpha_k\le \omega_2\cdot c_1c_2\cdots c_k$;
  \item[(C)]there exists $\omega_3>0$, such that $n_k\underline\alpha_k\ge \omega_3\cdot c_1c_2\cdots c_{k-1}$.
\end{itemize}

Then
\begin{equation}\label{311}
\dim_HE=\liminf\limits_{k \to\infty}  \frac{\log n_1n_2\cdots n_k}{-\log(\delta_k-L_{k+1}-R_{k+1})}.
\end{equation}
\end{thm}

\begin{thm}\label{thm2}
	    Let
	$E\in \mathcal{M}(I_{0},\left\{n_{k}\right\},\left\{c_{k}\right\})$ be a homogeneous Moran set which satisfies the following condition:
there exist two sequences of nonnegative real numbers $\{L_k\}_{k\ge1}$ and $\{R_k\}_{k\ge 1}$, such that
\begin{equation}
 \begin{aligned}
    \eta_{\sigma,0}=L_{k+1},\quad
    \eta_{\sigma,n_{k+1}}=R_{k+1} \nonumber
 \end{aligned}
 \end{equation}
 for any $k\ge 0$, $\sigma\in D_{k}$.

If $\dim_{H}E=1$, and for any $k\ge 1$, at least one of the following two conditions is satisfied:
\begin{itemize}
  \item[(A)]there exists $\omega_1>0$, such that $\bar\alpha_k\le \omega_1 \underline\alpha_k$;
  \item[(B)]there exists $\omega_2>0$, such that $\bar\alpha_k\le \omega_2\cdot c_1c_2\cdots c_k$.
\end{itemize}

Then we have $\operatorname{dim}_{H}f(E)=1$ for any \text{\rm1}-dimensional quasisymmetric mapping $f$, which implies that $E$ is a quasisymmetrically  Hausdorff-minimal set.

\end{thm}

\begin{lem}\label{l6}
	Let $E=E(I_{0},\left\{n_{k}\right\},\left\{c_{k}\right\})$ be a homogeneous Moran set which satisfies the conditions of Theorem {\rm\ref{thm2}}, and $E(I^*_{0},\{n^*_{k}\},\{c^*_{k}\})$ be the first reconstructed form of $E$. Then there is a sequence of closed sets  whose length is decreasing, denoted by $\{T_{m}\}_{m\ge 0}$, such that $E=\cap_{k\ge0}E_{k}=\cap_{k\ge0}E_{k}^{*}=\cap_{m\ge0}T_{m}$. Furthermore, $\{T_{m}\}_{m\ge 0}$ satisfies the following conditions:
\begin{enumerate}
  \item[\textup{(1)}] For any $m\ge0$, we have $T_{m}=\bigcup_{t=1}^{p_m} F_{t}$, where $p_m \in [1,+\infty)\cap \mathbb{Z}^{+}$,  $\{F_{t}\}_{1\le t\le p_m}$ is a sequence of close intervals, which are called the branches of $T_{m}$, and satisfying $\text{int}(F_{i_1})\cap \text{int}(F_{j_1})=\emptyset$ for any $1\le i_1<j_1\le p_m$. Denote $\mathcal{T}_{m}=\{A: A ~~~~\text{is a branch of}~~~~ T_{m}\}$;
  \item[\textup{(2)}] $\left\{E_{k}^{*}\right\}_{k\ge0}$ is a subsequence of $\left\{T_{m}\right\}_{m\ge 0}$, where $T_{m_{k}}=E_{k}^{*}$ for any $k\ge 0$;

  \item[\textup{(3)}] If $E$ satiefies the condition {\rm(A)} of Theorem {\rm\ref{thm2}}, then there exists $M \in \mathbb{Z}^{+}$ with $M>2\omega_1$ such that each branch of $T_{m-1}$ contains at most $M^{2}$ branches of $T_{m}$ for any $m\ge1$; if $E$ satiefies the condition {\rm(B)} of Theorem {\rm\ref{thm2}}, then there exists $M \in \mathbb{Z}^{+}$ with $M>2(\omega_2+1)$ such that each branch of $T_{m-1}$ contains at most $M^{2}$ branches of $T_{m}$ for any $m\ge1$, where $\omega_1$, $\omega_2$ are the constants in Theorem {\rm\ref{thm2}};
		\item[\textup{(4)}]If $E$ satiefies the condition {\rm(A)} of Theorem {\rm\ref{thm2}}, then $\max_{I\in \mathcal{T}_m}\left | I \right |\le 2 \omega_1\min_{I\in \mathcal{T}_m}\left | I \right |$, and if $E$ satiefies the condition {\rm(B)} of Theorem {\rm\ref{thm2}}, then
$\max_{I\in \mathcal{T}_m}\left | I \right |\le 2 (\omega_2+1)\min_{I\in \mathcal{T}_m}\left | I \right |$ for any $m\ge0$.
\end{enumerate}
\end{lem}

\begin{lem}\label{l7}
	Let $E=E(I_{0},\left\{n_{k}\right\},\left\{c_{k}\right\})$ be a homogeneous Moran set which satisfies the conditions of Theorem \ref{thm2}, $\left\{T_{m}\right\}_{m\ge0}$ be the sequences in Lemma \text{\rm\ref{l6}},
$l(T_{m})$ be the total length of all branches of $T_{m}$.
Then for any $k\ge 1$,
\begin{equation}
l(T_{m_{k}})=N^*_k\delta^*_k.
\end{equation}
Furthermore, if $E$ satiefies the condition {\rm(A)} of Theorem {\rm\ref{thm2}}, then for any $k\ge 1$ and $m_{k-1}<m<m_k$,
\begin{equation}
(1-\frac{2\omega_1}{M})N^*_{k-1}\delta^*_{k-1}\le l(T_m) \le N^*_{k-1}\delta^*_{k-1}.
\end{equation}
If $E$ satiefies the condition {\rm(B)} of Theorem {\rm\ref{thm2}}, then for any $k\ge 1$ and $m_{k-1}<m<m_k$,
\begin{equation}
(1-\frac{2(\omega_2+1)}{M})N^*_{k-1}\delta^*_{k-1}\le l(T_m) \le N^*_{k-1}\delta^*_{k-1}.
\end{equation}

\end{lem}

\begin{lem}\label{l8}
For any $m\ge0$, we have $\Theta_m\ge 1-(M^2+1)\beta_m$.
\end{lem}

\begin{lem}\label{l9}
Suppose $\{w_m\}_{m\ge 0}$ is a sequence of non-negative real numbers, and
\begin{equation*}
\lim_{m \to \infty} \frac{1}{m} \sum_{i=0}^{m-1} w_i=0.
\end{equation*}
Then we have
\begin{equation*}
\lim_{m \to \infty} \frac{V(m,\varepsilon )}{m} =1,
 \end{equation*}
for any $\varepsilon\in (0,1)$, where $V(m,\varepsilon )=\#(\{0\le i\le m-1: w_i<\varepsilon\})$$(\#$ denotes the cardinality$)$.
\end{lem}

\begin{lem}\label{l10}

Let $E=E(I_{0},\left\{n_{k}\right\},\left\{c_{k}\right\})$ be a homogeneous Moran set which satisfies the conditions of Theorem {\rm\ref{thm2}}, $\{T_{m}\}_{m\ge 0}$ be the sequences of Lemma {\rm\ref{l6}}.


If $\dim_{H}E=1$, then we have
\begin{enumerate}
\item[\textup{(1)}]$\lim_{m\to \infty}\frac{\log_{M}l(T_{m}) }{m}=0$, where $l(T_{m})$ is the total length of all branches of $T_{m}$;
\item[\textup{(2)}]$\lim_{m\to \infty}\frac{\sum_{j=0}^{m-1}\beta_{j}}{m}=0$;
\item[\textup{(3)}]$\lim_{m \to \infty} \frac{1}{m} \sum_{j=0}^{m-1}\log\Theta_j=0$;
\item[\textup{(4)}]there exists $\alpha\in (0,1)$, such that $\lim_{m\to \infty}  \frac{\# S(m,\alpha )}{m}=1$, where $S(m,\varepsilon)=\{1\le i\le m:\chi_i<\varepsilon\}$ for any $\varepsilon\in (0,1)$.
\end{enumerate}

\end{lem}

\begin{proof}
(1)
For any $m\ge 1$, if there exists $k\ge1$, such that $m=m_{k}$, then $l(T_m )=l(T_{m_k}) =l(E^*_k)=n^*_1n^*_2\cdots n^*_k\delta^*_k$. Since $\delta^*_kn^*_1n^*_2\cdots n^*_k=l(E^*_k)\le 1$, we have $\frac{\log_M(n^*_1n^*_2\cdots n^*_k)}{-\log_M\delta^*_k}\le 1.$ Notice that $\dim_{H}E=1$, by Theorem \ref{thm}, we have
\begin{equation}\begin{aligned}
1\ge\limsup\limits_{k \to\infty}  \frac{\log_M n^*_1n^*_2\cdots n^*_k}{-\log_M\delta^*_k}&\ge\liminf\limits_{k \to\infty}  \frac{\log_M n^*_1n^*_2\cdots n^*_k}{-\log_M\delta^*_k}\\& =\liminf\limits_{k \to\infty}  \frac{\log_M n_1n_2\cdots n_k}{-\log_M(\delta_k-L_{k+1}-R_{k+1})}\\& =\dim_{H}E=1.
\end{aligned}
\end{equation}
Which implies that \begin{equation}\label{101}
\lim_{k \to\infty}  \frac{\log_M n^*_1n^*_2\cdots n^*_k}{-\log_M\delta^*_k}=1.
\end{equation}

Since $\log_Mn^*_j\le i_j+1$, $m_k=i_1+i_2+\cdots+i_k$ and $m_k\ge k$ for any $k\ge 1$ and $1\le j \le k$, by (\ref{101}), we have
\begin{equation*}
\begin{aligned}
0\ge\lim_{k \to \infty} \frac{\log_M n^*_1n^*_2\cdots n^*_k\delta^*_k}{m_k} & = \lim_{k \to \infty}\frac{\log_M(n^*_1n^*_2\cdots n^*_k)}{m_k}\frac{\log_M(n^*_1n^*_2\cdots n^*_k)+\log_M\delta^*_k}{\log_M(n^*_1n^*_2\cdots n^*_k)}\\
&\ge \lim_{k \to \infty}2[1-( \frac{\log_M n^*_1n^*_2\cdots n^*_k}{-\log_M\delta^*_k})^{-1}]
=0.
\end{aligned}
\end{equation*}
Which implies that \begin{equation}\label{102}
\lim_{k \to \infty}\frac{\log_{M}l(T_{m_k}) }{m_k}=\lim_{k \to \infty} \frac{\log_M n^*_1n^*_2\cdots n^*_k\delta^*_k}{m_k}=0.
\end{equation}

If $E$ satiefies the condition {\rm(A)} of Theorem {\rm\ref{thm2}}, by Lemma \ref{l7}, we have $l(T_m) \ge(1-\frac{2\omega_1}{M})l(T_{m_{k-1}})$ for any $k\ge1 $ and $m_{k-1}\le m<m_k$. Notice that for any $\varepsilon>0$, there exists $N>0$, such that
$\frac{\log_{M} l(T_{m_k}) }{m_k}>-\frac{\varepsilon }{2} $ and $\frac{\log_{M}(1-\frac{2\omega_1}{M})^{-1} }{m_k}<\frac{\varepsilon }{2} $ for any $k\ge N$.
 Therefore if $m\ge m_{N}$, there is $h\ge N$ such that $m_h \le m < m_{h+1}$, then we have $$\frac{\log_M l(T_m)}{m}\ge\frac{\log_M l(T_{m_h})+\log_M(1-\frac{2\omega_1}{M})}{m_h}>-\varepsilon ,$$
which implies $\liminf\limits_{k \to\infty}   \frac{\log_M l(T_m ) }{m}=0$.
Since $l(T_m)\le 1$ for any $m\ge 0$,
$$\limsup\limits_{k \to\infty}  \frac{\log_M l(T_m) }{m}\le 0,$$ which implies that $$\lim_{m \to \infty} \frac{\log_M l(T_m) }{m}=0.$$






If $E$ satiefies the condition {\rm(B)} of Theorem {\rm\ref{thm2}}, by the similar proof (repalce $1-\frac{2\omega_1}{M}$ by $1-\frac{2(\omega_2+1)}{M}$), we have
$$\lim_{m \to \infty} \frac{\log_M l(T_m) }{m}=0.$$

(2) If $E$ satiefies the condition {\rm(A)} of Theorem {\rm\ref{thm2}}, for any $m_{k-1}\le m <m_k$, let $\kappa_m=\min\{\frac{\eta^*_{\sigma,l}}{|I|}:I\in\mathcal{T}_m, \sigma \in D_{k-1}, 1\le l\le n_k-1\}$. By the condition (A) of Theorem \ref{thm2} and (4) of Lemma \ref{l6}, we obtain $\beta_m\le 2\omega_1^2\kappa_m$. For any $0\le j\le m-1$ and $I\in T_j$, $\frac{T_{j+1}}{T_j}\le \frac{|I|-\kappa_j|I|}{|I|}=1-\kappa_j$, then  we have $l(T_m)\le \prod_{j=0}^{m-1}(1-\kappa_j)$. By the inequality $\log_M(1-x)\le -x$ for any $x\in[0,1)$, combining with (1), we have
$$0\ge \lim_{m\to \infty}-\frac{1}{m}\sum_{j=0}^{m-1}\kappa_j\ge \lim_{m\to \infty}\frac{1}{m}\sum_{j=0}^{m-1}\log(1-\kappa_j)\ge \lim_{m\to \infty}\frac{\log_{M}l(T_{m}) }{m}=0,$$
which implies that $$\lim_{m\to \infty}\frac{1}{m}\sum_{j=0}^{m-1}\kappa_j=0.$$
Then we have
$$0=2\omega^2\lim_{m\to \infty}\frac{1}{m}\sum_{j=0}^{m-1}\kappa_j\ge \lim_{m\to \infty}\frac{1}{m}\sum_{j=0}^{m-1}\beta_j\ge 0,$$
which implies that
$$\lim_{m\to \infty}\frac{1}{m}\sum_{j=0}^{m-1}\beta_j=0.$$

If $E$ satiefies the condition {\rm(B)} of Theorem {\rm\ref{thm2}}, by (1), we have
\begin{equation*}
\begin{aligned}
\lim_{k \to \infty} &\frac{1}{m_k} \log_{M}\prod_{j = 1}^{k}\frac{\delta^*_{j-1}-e^*_{j}-(L^*_j+R^*_j)}{\delta^*_{j-1}}\\
&= \lim_{k \to \infty}\frac{1}{m_k}\log_{M}\prod_{j = 1}^{k}\frac{n^*_j\delta^*_j}{\delta^*_{j-1}}\\
&= \lim_{k \to \infty}(\frac{1}{m_k}\log_{M}n^*_1n^*_2\cdots n^*_k\delta^*_k-\frac{1}{m_k}\log_{M}\delta^*_0)\\
&= \lim_{k \to \infty}(\frac{1}{m_k}\log_{M}l(S_{m_k}) -\frac{1}{m_k}\log_{M}\delta^*_0)\\
&=0.
\end{aligned}
\end{equation*}

By the inequality $\log_M(1-x)\le -x$ for any $x\in[0,1)$, we have $$0=\lim_{k \to \infty} \frac{1}{m_k} \log_{M}\prod_{j = 1}^{k}\frac{\delta^*_{j-1}-e^*_{j}-(L^*_j+R^*_j)}{\delta^*_{j-1}}\le\lim_{k \to \infty} -\frac{1}{m_k} \sum_{j = 1}^{k}\frac{e^*_{j}+(L^*_j+R^*_j)}{\delta^*_{j-1}}\le 0,$$
which implies that
$$\lim_{k \to \infty} \frac{1}{m_k} \sum_{j = 1}^{k}\frac{e^*_{j}+(L^*_j+R^*_j)}{\delta^*_{j-1}}=0.$$
Notice that for any $1\le j \le k$, $\bar\alpha^*_{j}=\bar\alpha_{j}+L^*_j+R^*_j\le e^*_{j}+L^*_j+R^*_j$, then
\begin{equation}\label{65}
\lim_{k \to \infty} \frac{1}{m_k} \sum_{j = 1}^{k}\frac{\bar{\alpha}^*_j}{\delta^*_{j-1}}=0.
\end{equation}

Next we estimate  $\beta_m$ for $m\ge0$. Let $k\in \mathbb{N} $ satisfying $m_{k-1}\le m < m_k$.
If $I\in T_{m_{k}-1}$, $I$ contains at least $2$ branches of $T_{m_k}$, notice that $\underline{\alpha}^*_k\ge L^*_k+R^*_k $,  therefore $\left | I \right | \ge 2\delta^*_k+(L^*_k+R^*_k)$.
If $I\in T_{m_{k}-2}$, $I$ contains at least $2M$ branches of $T_{m_k}$, therefore $\left | I \right | \ge 2M\delta^*_k+(2M-1)(L^*_k+R^*_k)$.
If $t\in \{1,2, \cdots, m_k-m_{k-1}\}$, then for any $I \in T_{m_k-t}$, $I$ contains at least $2M^{t-1}$ branches of $T_{m_k}$, therefore $\left | I \right | \ge 2M^{t-1}\delta^*_k+(2M^{t-1}-1)(L^*_k+R^*_k)$.
For any $L\in \mathcal{G}_m$, we have $\left | L \right | \le \bar{\alpha}^*_k,$
then for any $t\in \{1,2, \cdots, m_k-m_{k-1}\}$, we obtain
\begin{equation}
\begin{aligned}
\beta_{m_k-t}&\le \frac{\bar{\alpha}^*_k}{2M^{t-1}\delta^*_k+(2M^{t-1}-1)(L^*_k+R^*_k)}\\
&\le \frac{\bar{\alpha}^*_k}{2^{t}\delta^*_k+(2^{t}-1)(L^*_k+R^*_k)}\\
&\le \frac{\bar{\alpha}^*_k}{2^{t-1}(\delta^*_k+L^*_k+R^*_k)}.
\end{aligned}
\end{equation}
Therefore,
\begin{equation}\label{xxxx}
\begin{aligned}
\sum_{m = m_{k-1}}^{m_k-1} \beta_{m} = \sum_{t  = 1}^{i_k} \beta_{m_k-t}&\le\frac{\bar{\alpha}^*_k}{\delta^*_k+L^*_k+R^*_k}\sum_{t = 1}^{i_k}\frac{1}{2^{t-1}} \\
 &=\frac{\bar{\alpha}^*_k}{\delta^*_k+L^*_k+R^*_k}\sum_{t =0}^{i_k-1}\frac{1}{2^{t}}\\
&\le \frac{2\bar{\alpha}^*_k}{\delta^*_k+L^*_k+R^*_k}.
\end{aligned}
\end{equation}
Then we have
\begin{equation}\label{xs}
\frac{1}{m_k} \sum_{j=0}^{m_k-1} \beta_{j} \le \frac{2}{m_k}\sum_{j=1}^{k} \frac{\bar{\alpha}^*_j}{\delta^*_j+L^*_j+R^*_j}.
\end{equation}

For any $\varepsilon>0$, there exists $\delta>0$, such that
\begin{equation}\label{69}
0<\frac{1+\omega_2}{\log_M\frac{1}{(1+\omega_2)\delta}-1} <\frac{\varepsilon}{4}.
\end{equation}
If $j\ge1$ satisfying $\frac{\delta^*_j+L^*_j+R^*_j}{\delta^*_{j-1}} <\delta$.
Since $\bar{\alpha}_j\le \omega_2 \delta_j=\omega_2(\delta^*_j+L^*_j+R^*_j),$ we have
\begin{equation*}
\begin{aligned}
\delta^*_{j-1} & = e^*_{j}+L^*_j+R^*_j+n^*_j\delta^*_j\\
&=e_j+n^*_j(L^*_j+R^*_j)+n^*_j\delta^*_j\\
&\le (n^*_j-1)\bar{\alpha}_j+n^*_j(L^*_j+R^*_j+\delta^*_j)\\
&\le (n^*_j-1)\omega_2 (L^*_j+R^*_j+\delta^*_j)+n^*_j(L^*_j+R^*_j+\delta^*_j)\\
&\le n^*_j(1+\omega_2)(L^*_j+R^*_j+\delta^*_j),
\end{aligned}
\end{equation*}
which implies $(1+\omega_2)n^*_j\delta>(1+\omega_2)n^*_j\frac{L^*_j+R^*_j+\delta^*_j}{\delta^*_{j-1}} \ge 1.$
Notice that $i_j\ge \log_{M}n^*_j-1$, by (\ref{69}), we have
\begin{equation}\label{ss}
\frac{1+\omega_2}{i_j} <\frac{\varepsilon }{4}.
\end{equation}

By (\ref{65}), there exists $M_2>0$, such that for any $k\ge M_2$,
\begin{equation}\label{212}
\frac{1}{m_k} \sum_{j = 1}^{k}\frac{\bar{\alpha}^*_j}{\delta^*_{j-1}}<\frac{\varepsilon \delta}{4}.
\end{equation}
Notice that if $\bar{\alpha}_j\le \omega_2 \delta_j$, then
 \begin{align}\label{lwl}
\overline{\alpha } _{j}^{*}=\overline{\alpha}_j+L^*_j+R^*_j\le (\omega_2+1)\delta_j=(\omega_2+1)(\delta^*_{j}+L^*_j+R^*_j).
\end{align}
Therefore, if $k\ge M_2$, by (\ref{ss}),  (\ref{212}) and  (\ref{lwl}), we have
\begin{equation*}
\begin{aligned}
\frac{1}{m_k} \sum_{j = 1}^{k}\frac{\bar{\alpha}^*_j}{\delta^*_{j}+L^*_j+R^*_j}&\le \frac{1}{m_k}(\sum_{\substack{j=1\\\frac{\delta^*_j+L^*_j+R^*_j}{\delta^*_{j-1}}<\delta}}^{k}(1+\omega_2)+\sum_{\substack{j=1\\\frac{\delta^*_j+L^*_j+R^*_j}{\delta^*_{j-1}}\ge\delta}}^{k}\frac{\bar{\alpha}^*_j}{\delta^*_{j}+L^*_j+R^*_j} )\\
&\le \frac{1}{m_k}\sum_{j = 1}^{k}\frac{i_j\varepsilon }{4}+\frac{1}{m_k}\sum_{j = 1}^{k}\frac{\bar{\alpha}^*_j}{\delta^*_{j-1}}\cdot \frac{1}{\delta}\\
&<\frac{\varepsilon }{4}+\frac{\varepsilon }{4}=\frac{\varepsilon}{2},
\end{aligned}
\end{equation*}
which implies  $$\lim_{k \to \infty} \frac{1}{m_k} \sum_{j = 1}^{k}\frac{\bar{\alpha}^*_j}{\delta^*_{j}+L^*_j+R^*_j}=0.$$
Combing (\ref{xs}), we have
\begin{equation}\label{qwe}
\lim_{k \to \infty}\frac{1}{m_k} \sum_{j=0}^{m_k-1} \beta_{j}=0.
\end{equation}

For any  $k\ge 1$, if  $m_{k-1}<m<m_k$, then we have
\begin{equation}
\frac{1}{m} \sum_{j=0}^{m-1} \beta_{j}=\frac{1}{m} (\sum_{j=0}^{m_{k-1}-1} \beta_{j}+\sum_{j=m_{k-1}-1}^{m-1} \beta_{j})\le \frac{1}{m_{k-1}} \sum_{j=0}^{m_{k-1}-1} \beta_{j}+\frac{1}{m}\sum_{j=m_{k-1}-1}^{m-1}\beta_{j}.
\end{equation}
Notice that by (\ref{xxxx}) and (\ref{lwl}),
\begin{equation}\label{qas}
\sum_{j=m_{k-1}-1}^{m-1}\beta_{j}\le\frac{2\bar{\alpha}^*_k}{\delta^*_k+L^*_k+R^*_k}\le 2(1+\omega_2).
\end{equation}
Therefore, combing (\ref{qwe}) and (\ref{qas}), we obtain
\begin{equation*}
\lim_{m \to \infty} \frac{1}{m} \sum_{j=0}^{m-1} \beta_{j}=0.
\end{equation*}

(3) Fix $\varepsilon \in(0,\frac{1}{M^2+1} )$, such that $\log(1-(M^2+1)x)\ge -2(M^2+1)x$
for any $x\in[0,\varepsilon)$.
Then we have
\begin{equation*}
\begin{aligned}
0\ge \frac{1}{m} \sum_{\substack{j = 0\\\beta_j<\varepsilon }}^{m-1}\log(1-(M^2+1)\beta_j)\ge\frac{-2}{m} \sum_{\substack{j = 0\\\beta_j<\varepsilon }}^{m-1}(M^2+1)\beta_j&\ge
 -2(M^2+1)(\frac{1}{m}\sum_{j=0}^{m-1}\beta_j ).
\end{aligned}
\end{equation*}
By (2),
we have
\begin{equation*}
\lim_{m \to \infty} (\frac{1}{m} \sum_{\substack{j = 0\\\beta_j<\varepsilon }}^{m-1}\log(1-(M^2+1)\beta_j))=0,
\end{equation*}
which implies that
\begin{equation}\label{3a1}
\lim_{m \to \infty} ( \prod_{\substack{j = 0\\\beta_j<\varepsilon }}^{m-1}(1-(M^2+1)\beta_j))^{\frac{1}{m} }=1.
\end{equation}


If $E$ satiefies the condition {\rm(A)} of Theorem {\rm\ref{thm2}}, then each branch of $T_{m-1}$ contains at most $M^{2}$ branches of $T_{m}$ for any $m\ge1$, then we have $\frac{l(T_j) }{l(T_{j-1})}\le \min\{1,M^2\Lambda ^*(j)\}$ for any $1\le j \le m$. By (4) of Lemma \ref{l6}, we have $\Lambda^*(j)\le4\omega_1^2\Lambda_*(j)$ for any $j\ge 1$,
thus
\begin{equation}\label{3a2}
\prod_{j\in\Omega } (4M^2\omega_1^2\Lambda_*(j))\ge \prod_{j\in\Omega } (M^2\Lambda_*(j))\ge l(T_m)
\end{equation}
for any set $\Omega \subset \{1,2,\cdots,m\}$.


Let $H(m,\varepsilon)=\# (\{0\le j \le m-1:\beta_j<\varepsilon\}),$ combing (2) and Lemma \ref{l9}, we conclude that
\begin{equation}\label{3a3}
\lim_{m \to \infty} (1-\frac{H(m,\varepsilon )}{m}) =0.
\end{equation}

Notice that  for any $j\ge 0$, \begin{equation}\label{3a4}
\Theta_j=\min\{\frac{\sum_{i=1}^{N(I_j) }\left | I_{j,i} \right | }{\left | I_j \right | }: I_j\in T_j \}\ge \frac{\min_{I\in \mathcal{S}_{j+1} }\left | I \right | }{\max_{I\in \mathcal{S}_{j} }\left | I \right |}=\Lambda_*(j+1).
\end{equation}
Combing (\ref{3a2}), (\ref{3a4}) and Lemma \ref{l8}, we obtain


\begin{equation*}
\begin{aligned}
\lim_{m \to \infty} ( \prod_{j  = 0}^{m-1}\Theta_j)^{\frac{1}{m} } & =\lim_{m \to \infty} ( \prod_{\substack{j  = 0\\\beta_j<\varepsilon }}^{m-1}\Theta_j)^{\frac{1}{m} }( \prod_{\substack{j  = 0\\\beta_j\ge\varepsilon }}^{m-1}\Theta_j)^{\frac{1}{m} }\\
&\ge  \lim_{m \to \infty}(\prod_{\substack{j  = 0\\\beta_j<\varepsilon }}^{m-1}(1-(M^2+1)\beta_j))^{\frac{1}{m} }( \prod_{\substack{j  = 0\\\beta_j\ge\varepsilon }}^{m-1}\frac{1}{4M^2\omega_1^2} )^{\frac{1}{m} }l(T_m)^{\frac{1}{m} } \\
&=\lim_{m \to \infty}(\prod_{\substack{j  = 0\\\beta_j<\varepsilon }}^{m-1}(1-(M^2+1)\beta_j))^{\frac{1}{m} }(\frac{1}{4M^2\omega_1^2} )^{1-\frac{H(m,\varepsilon )}{m} }l(T_m)^{\frac{1}{m} }.
\end{aligned}
\end{equation*}

By (\ref{3a1}), we have $$\lim_{m \to \infty}(\prod_{\substack{j  = 0\\\beta_j<\varepsilon }}^{m-1}(1-(M^2+1)\beta_j))^{\frac{1}{m} }=1.$$
And by (\ref{3a3}) and (1), $$\lim_{m \to \infty}(\frac{1}{4M^2\omega_1^2} )^{1-\frac{H(m,\varepsilon )}{m} }l(T_m)^{\frac{1}{m} }=1.$$
Then $\lim_{m \to \infty} ( \prod_{j  = 0}^{m-1}\Theta_j)^{\frac{1}{m} }=1$, which implies
\begin{equation*}
\lim_{m \to \infty}\frac{1}{m} \sum_{j=0}^{m-1}\log \Theta_j=0 .
\end{equation*}




If $E$ satiefies the condition {\rm(B)} of Theorem {\rm\ref{thm2}}, by the similar proof(repalce $\omega_1$ by $\omega_2+1$), we also have
\begin{equation*}
\lim_{m \to \infty}\frac{1}{m} \sum_{j=0}^{m-1}\log \Theta_j=0 .
\end{equation*}




(4) If $E$ satiefies the condition {\rm(A)} of Theorem {\rm\ref{thm2}}, by (4) of Lemma \ref{l6}, for any $j\ge 1$, $J\in \mathcal{T}_j$ and $J^{\prime}\in \mathcal{T}_{j-1}$, we have
\begin{equation*}
\chi_j\le\frac{\max_{\hat{\hat{J}}\in T_j}\left | \hat{\hat{J}} \right | }{\min_{\hat{J}\in T_{j-1}}\left | \hat{J} \right |}\le\frac{2\omega_1\left | J \right | }{\frac{1}{2\omega_1} \left |J^{\prime} \right |}\le 4\omega_1^2\frac{\left | J \right |}{\left |J^{\prime} \right |}.
\end{equation*}


Taking $J^0\in T_j$ which satisfies $\chi_j=\frac{\left | J^0 \right |}{\left |X_a(J^0) \right |} $.
Since $X_a(J^0)$  contains at least $2$ branches of $T_j$, we have
\begin{equation*}
1>\chi_j+\frac{\left | J^* \right |}{\left |X_a(J^0) \right |}\ge \chi_j+\frac{\chi_j}{4\omega_1^2}=(\frac{4\omega_1^2+1}{4\omega_1^2} )\chi_j,
\end{equation*}
where $J^*\in \mathcal{T}_j$, $J^*\subset X_a(J^0)$ and $J^*\neq J$. Let $\alpha \in (\frac{4\omega_1^2}{4\omega_1^2+1},1)$, then
\begin{equation*}
\lim_{m\to \infty} \frac{\#S(m,\alpha )}{m} =\lim_{m\to \infty}\frac{\#\{1\le i\le m:\chi_i<\alpha\}}{m} =1.
\end{equation*}


If $E$ satiefies the condition {\rm(B)} of Theorem {\rm\ref{thm2}},
 by the similar proof(repalce $\omega_1$ by $\omega_2+1$), we also have
\begin{equation*}
\lim_{m\to \infty} \frac{\#S(m,\alpha )}{m} =\lim_{m\to \infty} \frac{\#\{1\le i\le m:\chi_i<\alpha\}}{m} =1.
\end{equation*}





\end{proof}

\end{document}
